% Part: proof-theory
% Chapter: natural-deduction
% Section: translation-N2i

\documentclass[../../../include/open-logic-section]{subfiles}

\begin{document}

\olfileid[pt]{pt}{nat}{tni}

\olsection{Tradução de \Log{N2} para \Log{G2}}

\begin{prop}\ollabel{prop:N2-to-G2} Se $\Log{N2c}\ (\Log{N2i}) \Proves
\Gamma \Sequent !A$, então $\Log{G2c}\ (\Log{N2i}) + \Cut \Proves \Gamma'
\Sequent \Delta$, em que $\Gamma'$ é o multiconjunto de !!{formula}s
obtido de $\Gamma$ removendo os rótulos, e $\Delta = \{!A\}$ se $!A$
não é~$\lfalse$, enquanto $\Delta = \emptyset$ se é.
\end{prop}

\begin{proof}
A demonstração é novamente por indução sobre a altura da
!!{proof}~$\delta$ de $\Gamma \Sequent !A$. Se $\delta$ é o axioma
$x:!A \Sequent !A$, então $\pi$ é o axioma correspondente $!A \Sequent !A$
ou $\lfalse \Sequent \ $, caso $!A \ident \lfalse$.

Se $\delta$ termina em uma regra de inferência, distinguimos os casos:
\begin{enumerate}
  \item Se $R$ é \Intro{\lif}, com fórmula principal $!B \lif !C$, e
  $x:!B$ está presente no contexto da premissa, mas não no da
  conclusão, então $\delta$ termina em
  \[
  \AxiomC{}
  \RightLabel{$\delta_1$}
  \Deduce$x:!B, \Gamma \fCenter !C$
  \RightLabel{\Intro{\lif}}
  \UnaryInf$\Gamma \fCenter !B \lif !C$
  \DisplayProof
  \]
  A hipótese indutiva fornece !!a{proof}~$\pi_1$ de $!B, \Gamma'
  \Sequent !C$ (ou de $!B, \Gamma' \Sequent \ $, se $!C \ident
  \lfalse$), da qual obtemos $\pi$ aplicando~\RightR{\lif} (depois de
  aplicar \RightR{\Weakening}, caso $!C \ident \lfalse$).
  \item Se $R$ é \Intro{\lif}, mas $x:!B$ não está presente no
  contexto da premissa, então $\delta$ termina em
  \[
  \AxiomC{}
  \RightLabel{$\delta_1$}
  \Deduce$\Gamma \fCenter !C$
  \RightLabel{\Intro{\lif}}
  \UnaryInf$\Gamma \fCenter !B \lif !C$
  \DisplayProof
  \]
  A hipótese indutiva fornece !!a{proof}~$\pi_1$ de $\Gamma'
  \Sequent !C$, da qual obtemos $\pi$ aplicando
  \LeftR{\Weakening} com~$!B$ e, depois,~\RightR{\lif}.
  \item Se $R$ é \Elim{\lif}, então $\delta$ termina em
  \[
  \AxiomC{}
  \RightLabel{$\delta_1$}
  \Deduce$\Gamma_1 \fCenter !B \lif !A$
  \AxiomC{}
  \RightLabel{$\delta_2$}
  \Deduce$\Gamma_2 \fCenter !B$
  \RightLabel{\Elim{\lif}}
  \BinaryInf$\Gamma_1, \Gamma_2 \fCenter !A$
  \DisplayProof
  \]
  em que $\Gamma = \Gamma_1 \cup \Gamma_2$. A hipótese indutiva
  fornece !!{proof}s $\pi_1$ de $\Gamma_1' \Sequent !B \lif !A$ e
  $\pi_2$ de $\Gamma_2' \Sequent !B$. Obtemos a !!{proof}~$\pi$ como
  \[
  \AxiomC{}
  \RightLabel{$\pi_1$}
  \Deduce$\Gamma_1' \fCenter !B \lif !A$
  \AxiomC{}
  \RightLabel{$\pi_2$}
  \Deduce$\Gamma_2' \fCenter !B$
  \Axiom$!A \fCenter !A$
  \RightLabel{\LeftR{\lif}}
  \BinaryInf$!B \lif !A, \Gamma_2' \fCenter !A$
  \RightLabel{\Cut}
  \BinaryInf$\Gamma_1', \Gamma_2' \fCenter !A$
  \DisplayProof
  \]
  O multiconjunto $\Gamma_1' \cup \Gamma_2'$ pode não ser igual a
  $\Gamma' = (\Gamma_1 \cup \Gamma_2)'$: por exemplo, se tanto
  $\Gamma_1$ quanto $\Gamma_2$ contêm $x:!C$, então $\Gamma_1' \cup
  \Gamma_2'$ contém duas cópias de $!C$, mas $(\Gamma_1 \cup
  \Gamma_2)'$ contém apenas uma. Podemos corrigir isso acrescentando
  inferências \LeftR{\Contraction} apropriadas.

  No caso em que $!A \ident \lfalse$, a hipótese indutiva aplicada a
  $\delta_2$ fornece !!a{proof} de $\Gamma_2' \Sequent \ $, da qual
  podemos obter a !!{proof}~$\pi$ acrescentando inferências
  \LeftR{\Weakening} apropriadas e uma \RightR{\Weakening} com~$!A$.
\item Se $R$ é \FalseCl, então $\delta$ termina em
  \[
  \AxiomC{}
  \RightLabel{$\delta_1$}
  \Deduce$x: \lnot !A, \Gamma \fCenter \lfalse$
  \RightLabel{\FalseCl}
  \UnaryInf$\Gamma \fCenter !A$
  \DisplayProof
  \]
  Pela hipótese indutiva, existe !!a{proof} $\pi_1$ de $\lnot !A,
  \Gamma' \Sequent \ $. Tomamos como !!{proof}~$\pi$
  \[
  \Axiom$!A \fCenter !A$
  \RightLabel{\RightR{\lnot}}
  \UnaryInf$\fCenter !A, \lnot !A$
  \AxiomC{}
  \RightLabel{$\pi_1$}
  \Deduce$x: \lnot !A, \Gamma' \fCenter $
  \RightLabel{\Cut}
  \BinaryInf$\Gamma' \fCenter !A$
  \DisplayProof
  \]
\end{enumerate}
Observe que apenas a tradução de \FalseCl{} resulta em
!!a{proof}~$\pi$ com mais de uma !!{formula} à direita. Isto é, a
tradução de uma !!{proof} em~\Log{N2i} é uma \Log{G2i}-!!{proof}.
\end{proof}

\end{document}
